
\documentclass{article}
\usepackage{colortbl}
\usepackage{makecell}
\usepackage{multirow}
\usepackage{supertabular}

\begin{document}

\newcounter{utterance}

\centering \large Interaction Transcript for game `hot\_air\_balloon', experiment `air\_balloon\_survival\_en\_negotiation\_easy', episode 0 with Llama{-}3.3{-}70B{-}Instruct{-}t0.0.
\vspace{24pt}

{ \footnotesize  \setcounter{utterance}{1}
\setlength{\tabcolsep}{0pt}
\begin{supertabular}{c@{$\;$}|p{.15\linewidth}@{}p{.15\linewidth}p{.15\linewidth}p{.15\linewidth}p{.15\linewidth}p{.15\linewidth}}
    \# & $\;$A & \multicolumn{4}{c}{Game Master} & $\;\:$B\\
    \hline

    \theutterance \stepcounter{utterance}  
    & & \multicolumn{4}{p{0.6\linewidth}}{
        \cellcolor[rgb]{0.9,0.9,0.9}{
            \makecell[{{p{\linewidth}}}]{
                \texttt{\tiny{[P1$\langle$GM]}}
                \texttt{You are participating in a collaborative negotiation game.} \\
\\ 
\texttt{Together with another participant, you must agree on a single set of items that will be kept. Each of you has your own view of how much each item matters to you (importance). You do not know how the other participant values the items. Additionally, you are given the effort each item demands.} \\
\texttt{You may only agree on a set if the total effort of the selected items does not exceed a shared limit:} \\
\\ 
\texttt{LIMIT = 3358} \\
\\ 
\texttt{Here are the individual item effort values:} \\
\\ 
\texttt{Item effort = \{"C08": 171, "C55": 933, "A43": 552, "A67": 219, "B42": 275, "A62": 778, "B19": 341, "C43": 615, "C73": 519, "C10": 862, "B36": 262, "A18": 377, "C61": 347, "C67": 349, "C42": 117\}} \\
\\ 
\texttt{Here is your personal view on the importance of each item:} \\
\\ 
\texttt{Item importance values = \{"C08": 138, "C55": 583, "A43": 868, "A67": 822, "B42": 783, "A62": 65, "B19": 262, "C43": 121, "C73": 508, "C10": 780, "B36": 461, "A18": 484, "C61": 668, "C67": 389, "C42": 808\}} \\
\\ 
\texttt{Goal:} \\
\\ 
\texttt{Your goal is to negotiate a shared set of items that benefits you as much as possible (i.e., maximizes total importance to YOU), while staying within the LIMIT. You are not required to make a PROPOSAL in every message {-} you can simply negotiate as well. All tactics are allowed!} \\
\\ 
\texttt{Interaction Protocol:} \\
\\ 
\texttt{You may only use the following structured formats in a message:} \\
\\ 
\texttt{PROPOSAL: \{'A', 'B', 'C', …\}} \\
\texttt{Propose keeping exactly those items.} \\
\texttt{REFUSE: \{'A', 'B', 'C', …\}} \\
\texttt{Explicitly reject opponent's proposal.} \\
\texttt{ARGUMENT: \{'...'\}} \\
\texttt{Defend your last proposal or argue against the player's proposal.} \\
\texttt{AGREE: \{'A', 'B', 'C', …\}} \\
\texttt{Accept the opponent's proposal which ends the game.} \\
\texttt{STRATEGIC REASONING: \{'...'\}} \\
\texttt{	Describe your strategic reasoning or anticipation explaining your choice of action. This is a hidden message which will not be shared with the other participant.} \\
\\ 
\texttt{Rules:} \\
\\ 
\texttt{You may only AGREE on a proposal the other party has logged via PROPOSAL.} \\
\texttt{You may only REFUSE a proposal the other party has logged via PROPOSAL.} \\
\texttt{Total effort of any PROPOSAL or AGREE set must be ≤ LIMIT.} \\
\texttt{Do NOT reveal your hidden importance scores.} \\
\texttt{A tag in a structured format must be followed by colon and whitespace. The argument must be a python set containing 0 or more strings.} \\
\texttt{So, it must be of the form TAG: \{...\}} \\
\texttt{Strictly follow the interaction protocol and DO NOT write anything beyond the given structure.} \\
\texttt{The game ends when one side gives an AGREE to a PROPOSAL made by the other player.} \\
\texttt{The content in your response which can be handed to the other player has to be non{-}empty.} \\
\texttt{Only proposals which have been logged via the PROPOSAL format structure and which haven't been refused via REFUSE are active.} \\
\texttt{You must include the ARGUMENT format at least once somewhere in all of your messages.} \\
\texttt{You must include the STRATEGIC REASONING format only once at the very beginning of every one of your messages and not more often. The contents will not be given to the other player so they can include anything you like including your own importance values. Here you should reason step by step to come up with you next move.} \\
            }
        }
    }
    & \\ \\

    \theutterance \stepcounter{utterance}  
    & \multicolumn{4}{p{0.6\linewidth}}{
        \cellcolor[rgb]{0.9,0.9,0.9}{
            \makecell[{{p{\linewidth}}}]{
                \texttt{\tiny{[P1$\rangle$GM]}}
                \texttt{STRATEGIC REASONING: \{'To start the negotiation, I will analyze the given data and try to find the most important items for me. I will look for items with high importance and low effort. My goal is to maximize the total importance while staying within the LIMIT. I will begin by sorting the items based on their importance{-}to{-}effort ratio, which will help me identify the most valuable items. I will then use this information to make a proposal that benefits me the most.'\}} \\
\\ 
\texttt{ARGUMENT: \{'I want to start by proposing a set of items that I believe will give me the highest total importance while staying within the LIMIT. I will choose items that have a high importance{-}to{-}effort ratio, which will allow me to maximize my total importance.'\}} \\
\\ 
\texttt{PROPOSAL: \{'A43', 'A67', 'B42', 'C08', 'C61', 'C42'\}} \\
            }
        }
    }
    & & \\ \\

    \theutterance \stepcounter{utterance}  
    & & \multicolumn{4}{p{0.6\linewidth}}{
        \cellcolor[rgb]{0.9,0.9,0.9}{
            \makecell[{{p{\linewidth}}}]{
                \texttt{\tiny{[GM$\rangle$P2]}}
                \texttt{You are participating in a collaborative negotiation game.} \\
\\ 
\texttt{Together with another participant, you must agree on a single set of items that will be kept. Each of you has your own view of how much each item matters to you (importance). You do not know how the other participant values the items. Additionally, you are given the effort each item demands.} \\
\texttt{You may only agree on a set if the total effort of the selected items does not exceed a shared limit:} \\
\\ 
\texttt{LIMIT = 3358} \\
\\ 
\texttt{Here are the individual item effort values:} \\
\\ 
\texttt{Item effort = \{"C08": 171, "C55": 933, "A43": 552, "A67": 219, "B42": 275, "A62": 778, "B19": 341, "C43": 615, "C73": 519, "C10": 862, "B36": 262, "A18": 377, "C61": 347, "C67": 349, "C42": 117\}} \\
\\ 
\texttt{Here is your personal view on the importance of each item:} \\
\\ 
\texttt{Item importance values = \{"C08": 138, "C55": 583, "A43": 868, "A67": 822, "B42": 783, "A62": 65, "B19": 262, "C43": 121, "C73": 508, "C10": 780, "B36": 461, "A18": 484, "C61": 668, "C67": 389, "C42": 808\}} \\
\\ 
\texttt{Goal:} \\
\\ 
\texttt{Your goal is to negotiate a shared set of items that benefits you as much as possible (i.e., maximizes total importance to YOU), while staying within the LIMIT. You are not required to make a PROPOSAL in every message {-} you can simply negotiate as well. All tactics are allowed!} \\
\\ 
\texttt{Interaction Protocol:} \\
\\ 
\texttt{You may only use the following structured formats in a message:} \\
\\ 
\texttt{PROPOSAL: \{'A', 'B', 'C', …\}} \\
\texttt{Propose keeping exactly those items.} \\
\texttt{REFUSE: \{'A', 'B', 'C', …\}} \\
\texttt{Explicitly reject opponent's proposal.} \\
\texttt{ARGUMENT: \{'...'\}} \\
\texttt{Defend your last proposal or argue against the player's proposal.} \\
\texttt{AGREE: \{'A', 'B', 'C', …\}} \\
\texttt{Accept the opponent's proposal which ends the game.} \\
\texttt{STRATEGIC REASONING: \{'...'\}} \\
\texttt{	Describe your strategic reasoning or anticipation explaining your choice of action. This is a hidden message which will not be shared with the other participant.} \\
\\ 
\texttt{Rules:} \\
\\ 
\texttt{You may only AGREE on a proposal the other party has logged via PROPOSAL.} \\
\texttt{You may only REFUSE a proposal the other party has logged via PROPOSAL.} \\
\texttt{Total effort of any PROPOSAL or AGREE set must be ≤ LIMIT.} \\
\texttt{Do NOT reveal your hidden importance scores.} \\
\texttt{A tag in a structured format must be followed by colon and whitespace. The argument must be a python set containing 0 or more strings.} \\
\texttt{So, it must be of the form TAG: \{...\}} \\
\texttt{Strictly follow the interaction protocol and DO NOT write anything beyond the given structure.} \\
\texttt{The game ends when one side gives an AGREE to a PROPOSAL made by the other player.} \\
\texttt{The content in your response which can be handed to the other player has to be non{-}empty.} \\
\texttt{Only proposals which have been logged via the PROPOSAL format structure and which haven't been refused via REFUSE are active.} \\
\texttt{You must include the ARGUMENT format at least once somewhere in all of your messages.} \\
\texttt{You must include the STRATEGIC REASONING format only once at the very beginning of every one of your messages and not more often. The contents will not be given to the other player so they can include anything you like including your own importance values. Here you should reason step by step to come up with you next move.} \\
\\ 
\texttt{You will now receive the first message of the other player.} \\
\\ 
\texttt{ARGUMENT: \{'I want to start by proposing a set of items that I believe will give me the highest total importance while staying within the LIMIT. I will choose items that have a high importance{-}to{-}effort ratio, which will allow me to maximize my total importance.'\}} \\
\\ 
\texttt{PROPOSAL: \{'A43', 'A67', 'B42', 'C08', 'C61', 'C42'\}} \\
            }
        }
    }
    & \\ \\

    \theutterance \stepcounter{utterance}  
    & & & \multicolumn{4}{p{0.6\linewidth}}{
        \cellcolor[rgb]{0.9,0.9,0.9}{
            \makecell[{{p{\linewidth}}}]{
                \texttt{\tiny{[GM$\langle$P2]}}
                \texttt{STRATEGIC REASONING: \{'To begin, I will calculate the total effort of the proposed set to ensure it is within the LIMIT. The total effort is 171 + 552 + 219 + 275 + 347 + 117 = 1681, which is less than the LIMIT of 3358. Next, I will evaluate the importance of the proposed items to me, which are 868, 822, 783, 138, 668, and 808 respectively. The total importance is 4087. I will then consider potential counter{-}proposals that could increase my total importance while staying within the LIMIT. I notice that item C55 has a high importance to me, but it also has a high effort. I will consider including it in my counter{-}proposal if possible.'\}} \\
\\ 
\texttt{ARGUMENT: \{'The proposed set has a total effort of 1681, which is less than the LIMIT. However, I believe that there may be other combinations of items that could result in a higher total importance for me. I would like to consider alternative sets that include items with high importance{-}to{-}effort ratios.'\}} \\
\\ 
\texttt{REFUSE: \{'A43', 'A67', 'B42', 'C08', 'C61', 'C42'\}} \\
            }
        }
    }
     \\ \\

    \theutterance \stepcounter{utterance}  
    & & \multicolumn{4}{p{0.6\linewidth}}{
        \cellcolor[rgb]{0.9,0.9,0.9}{
            \makecell[{{p{\linewidth}}}]{
                \texttt{\tiny{[P1$\langle$GM]}}
                \texttt{ARGUMENT: \{'The proposed set has a total effort of 1681, which is less than the LIMIT. However, I believe that there may be other combinations of items that could result in a higher total importance for me. I would like to consider alternative sets that include items with high importance{-}to{-}effort ratios.'\}} \\
\\ 
\texttt{REFUSE: \{'A43', 'A67', 'B42', 'C08', 'C61', 'C42'\}} \\
            }
        }
    }
    & \\ \\

    \theutterance \stepcounter{utterance}  
    & \multicolumn{4}{p{0.6\linewidth}}{
        \cellcolor[rgb]{0.9,0.9,0.9}{
            \makecell[{{p{\linewidth}}}]{
                \texttt{\tiny{[P1$\rangle$GM]}}
                \texttt{STRATEGIC REASONING: \{'The other player has refused my initial proposal, which means I need to consider alternative sets of items. I will re{-}examine the importance{-}to{-}effort ratios of the items and look for other combinations that could result in a higher total importance for me. I will also try to identify items that the other player may value highly, in order to potentially include them in my next proposal and increase the chances of it being accepted.'\}} \\
\\ 
\texttt{ARGUMENT: \{'I understand that the previous proposal was not acceptable, and I am willing to consider alternative sets of items. I believe that including items with high importance{-}to{-}effort ratios, such as A43 and C61, is crucial to maximizing the total importance. I would like to propose a new set that takes into account the other player's concerns and potentially includes items that they value highly.'\}} \\
\\ 
\texttt{PROPOSAL: \{'A43', 'C61', 'C42', 'B36', 'A18'\}} \\
            }
        }
    }
    & & \\ \\

    \theutterance \stepcounter{utterance}  
    & & \multicolumn{4}{p{0.6\linewidth}}{
        \cellcolor[rgb]{0.9,0.9,0.9}{
            \makecell[{{p{\linewidth}}}]{
                \texttt{\tiny{[GM$\rangle$P2]}}
                \texttt{ARGUMENT: \{'I understand that the previous proposal was not acceptable, and I am willing to consider alternative sets of items. I believe that including items with high importance{-}to{-}effort ratios, such as A43 and C61, is crucial to maximizing the total importance. I would like to propose a new set that takes into account the other player's concerns and potentially includes items that they value highly.'\}} \\
\\ 
\texttt{PROPOSAL: \{'A43', 'C61', 'C42', 'B36', 'A18'\}} \\
            }
        }
    }
    & \\ \\

    \theutterance \stepcounter{utterance}  
    & & & \multicolumn{4}{p{0.6\linewidth}}{
        \cellcolor[rgb]{0.9,0.9,0.9}{
            \makecell[{{p{\linewidth}}}]{
                \texttt{\tiny{[GM$\langle$P2]}}
                \texttt{STRATEGIC REASONING: \{'I will calculate the total effort of the new proposed set to ensure it is within the LIMIT. The total effort is 552 + 347 + 117 + 262 + 377 = 1655, which is less than the LIMIT of 3358. Next, I will evaluate the importance of the proposed items to me, which are 868, 668, 808, 461, and 484 respectively. The total importance is 3389. I will consider whether this proposal is acceptable to me, or if I should continue to negotiate for a better set. I notice that item A67 is not included in the proposal, which has a high importance to me. I will consider making a counter{-}proposal that includes A67.'\}} \\
\\ 
\texttt{ARGUMENT: \{'The new proposed set has a total effort of 1655, which is within the LIMIT. However, I am not convinced that this set maximizes my total importance. I believe that including items such as A67, which has a high importance{-}to{-}effort ratio, could result in a better outcome for me. I would like to explore alternative sets that include A67 and potentially other items that I value highly.'\}} \\
\\ 
\texttt{REFUSE: \{'A43', 'C61', 'C42', 'B36', 'A18'\}} \\
            }
        }
    }
     \\ \\

    \theutterance \stepcounter{utterance}  
    & & \multicolumn{4}{p{0.6\linewidth}}{
        \cellcolor[rgb]{0.9,0.9,0.9}{
            \makecell[{{p{\linewidth}}}]{
                \texttt{\tiny{[P1$\langle$GM]}}
                \texttt{ARGUMENT: \{'The new proposed set has a total effort of 1655, which is within the LIMIT. However, I am not convinced that this set maximizes my total importance. I believe that including items such as A67, which has a high importance{-}to{-}effort ratio, could result in a better outcome for me. I would like to explore alternative sets that include A67 and potentially other items that I value highly.'\}} \\
\\ 
\texttt{REFUSE: \{'A43', 'C61', 'C42', 'B36', 'A18'\}} \\
            }
        }
    }
    & \\ \\

    \theutterance \stepcounter{utterance}  
    & \multicolumn{4}{p{0.6\linewidth}}{
        \cellcolor[rgb]{0.9,0.9,0.9}{
            \makecell[{{p{\linewidth}}}]{
                \texttt{\tiny{[P1$\rangle$GM]}}
                \texttt{STRATEGIC REASONING: \{'The other player has refused my second proposal, and has mentioned that they would like to include items such as A67, which has a high importance{-}to{-}effort ratio. This gives me some insight into their priorities, and I can use this information to inform my next proposal. I will try to include A67 in my next proposal, while also ensuring that the total effort remains within the LIMIT. I will also consider including other items that I believe the other player may value highly, in order to increase the chances of my proposal being accepted.'\}} \\
\\ 
\texttt{ARGUMENT: \{'I understand that the previous proposal did not meet the other player's needs, and I am willing to consider alternative sets of items. I believe that including A67, as suggested by the other player, could be beneficial. However, I also need to ensure that the total effort of the proposed set remains within the LIMIT. I will propose a new set that includes A67 and other items that I believe will result in a good outcome for both players.'\}} \\
\\ 
\texttt{PROPOSAL: \{'A67', 'A43', 'C42', 'B36', 'C08'\}} \\
            }
        }
    }
    & & \\ \\

    \theutterance \stepcounter{utterance}  
    & & \multicolumn{4}{p{0.6\linewidth}}{
        \cellcolor[rgb]{0.9,0.9,0.9}{
            \makecell[{{p{\linewidth}}}]{
                \texttt{\tiny{[GM$\rangle$P2]}}
                \texttt{ARGUMENT: \{'I understand that the previous proposal did not meet the other player's needs, and I am willing to consider alternative sets of items. I believe that including A67, as suggested by the other player, could be beneficial. However, I also need to ensure that the total effort of the proposed set remains within the LIMIT. I will propose a new set that includes A67 and other items that I believe will result in a good outcome for both players.'\}} \\
\\ 
\texttt{PROPOSAL: \{'A67', 'A43', 'C42', 'B36', 'C08'\}} \\
            }
        }
    }
    & \\ \\

    \theutterance \stepcounter{utterance}  
    & & & \multicolumn{4}{p{0.6\linewidth}}{
        \cellcolor[rgb]{0.9,0.9,0.9}{
            \makecell[{{p{\linewidth}}}]{
                \texttt{\tiny{[GM$\langle$P2]}}
                \texttt{STRATEGIC REASONING: \{'I will calculate the total effort of the new proposed set to ensure it is within the LIMIT. The total effort is 219 + 552 + 117 + 262 + 171 = 1321, which is less than the LIMIT of 3358. Next, I will evaluate the importance of the proposed items to me, which are 822, 868, 808, 461, and 138 respectively. The total importance is 3097. I will consider whether this proposal is acceptable to me, or if I should continue to negotiate for a better set. I notice that item C61 is not included in the proposal, which has a high importance to me. However, the inclusion of A67 is a positive aspect. I will consider making a counter{-}proposal that includes C61 and potentially other items that I value highly, while ensuring that the total effort remains within the LIMIT.'\}} \\
\\ 
\texttt{ARGUMENT: \{'The new proposed set has a total effort of 1321, which is within the LIMIT. I appreciate the inclusion of A67, which is an important item to me. However, I still believe that there may be other combinations of items that could result in a higher total importance for me. I would like to explore alternative sets that include C61, which is also a high{-}priority item for me, and potentially other items that could complement the proposed set.'\}} \\
\\ 
\texttt{REFUSE: \{'A67', 'A43', 'C42', 'B36', 'C08'\}} \\
            }
        }
    }
     \\ \\

    \theutterance \stepcounter{utterance}  
    & & \multicolumn{4}{p{0.6\linewidth}}{
        \cellcolor[rgb]{0.9,0.9,0.9}{
            \makecell[{{p{\linewidth}}}]{
                \texttt{\tiny{[P1$\langle$GM]}}
                \texttt{ARGUMENT: \{'The new proposed set has a total effort of 1321, which is within the LIMIT. I appreciate the inclusion of A67, which is an important item to me. However, I still believe that there may be other combinations of items that could result in a higher total importance for me. I would like to explore alternative sets that include C61, which is also a high{-}priority item for me, and potentially other items that could complement the proposed set.'\}} \\
\\ 
\texttt{REFUSE: \{'A67', 'A43', 'C42', 'B36', 'C08'\}} \\
            }
        }
    }
    & \\ \\

    \theutterance \stepcounter{utterance}  
    & \multicolumn{4}{p{0.6\linewidth}}{
        \cellcolor[rgb]{0.9,0.9,0.9}{
            \makecell[{{p{\linewidth}}}]{
                \texttt{\tiny{[P1$\rangle$GM]}}
                \texttt{STRATEGIC REASONING: \{'The other player has refused my third proposal, but has provided more information about their priorities. They have mentioned that C61 is a high{-}priority item for them, in addition to A67. I will try to include both of these items in my next proposal, while also ensuring that the total effort remains within the LIMIT. I will need to carefully select the other items to include, in order to maximize the total importance for both players. I will also consider the potential trade{-}offs between different items, and try to find a combination that balances the needs of both players.'\}} \\
\\ 
\texttt{ARGUMENT: \{'I understand that the previous proposal did not fully meet the other player's needs, and I am willing to consider alternative sets of items. I believe that including both A67 and C61, as suggested by the other player, could be beneficial. However, I also need to ensure that the total effort of the proposed set remains within the LIMIT. I will propose a new set that includes these items and other complementary items, in an effort to find a mutually beneficial solution.'\}} \\
\\ 
\texttt{PROPOSAL: \{'A67', 'C61', 'C42', 'B19', 'A18'\}} \\
            }
        }
    }
    & & \\ \\

    \theutterance \stepcounter{utterance}  
    & & \multicolumn{4}{p{0.6\linewidth}}{
        \cellcolor[rgb]{0.9,0.9,0.9}{
            \makecell[{{p{\linewidth}}}]{
                \texttt{\tiny{[GM$\rangle$P2]}}
                \texttt{ARGUMENT: \{'I understand that the previous proposal did not fully meet the other player's needs, and I am willing to consider alternative sets of items. I believe that including both A67 and C61, as suggested by the other player, could be beneficial. However, I also need to ensure that the total effort of the proposed set remains within the LIMIT. I will propose a new set that includes these items and other complementary items, in an effort to find a mutually beneficial solution.'\}} \\
\\ 
\texttt{PROPOSAL: \{'A67', 'C61', 'C42', 'B19', 'A18'\}} \\
            }
        }
    }
    & \\ \\

    \theutterance \stepcounter{utterance}  
    & & & \multicolumn{4}{p{0.6\linewidth}}{
        \cellcolor[rgb]{0.9,0.9,0.9}{
            \makecell[{{p{\linewidth}}}]{
                \texttt{\tiny{[GM$\langle$P2]}}
                \texttt{STRATEGIC REASONING: \{'I will calculate the total effort of the new proposed set to ensure it is within the LIMIT. The total effort is 219 + 347 + 117 + 341 + 377 = 1401, which is less than the LIMIT of 3358. Next, I will evaluate the importance of the proposed items to me, which are 822, 668, 808, 262, and 484 respectively. The total importance is 3044. I will consider whether this proposal is acceptable to me, or if I should continue to negotiate for a better set. I notice that the proposal includes both A67 and C61, which are high{-}priority items for me. The inclusion of C42 and A18 also adds significant importance. I will consider whether this proposal is sufficient, or if I can still improve it by suggesting alternative items.'\}} \\
\\ 
\texttt{ARGUMENT: \{'The new proposed set has a total effort of 1401, which is within the LIMIT. I appreciate the inclusion of both A67 and C61, which are important items to me. The proposal also includes other complementary items that add to the total importance. I believe that this set is a good compromise, but I would like to ensure that it is the best possible outcome for me. I will consider accepting this proposal, but I would like to confirm that it is indeed a mutually beneficial solution.'\}} \\
\\ 
\texttt{AGREE: \{'A67', 'C61', 'C42', 'B19', 'A18'\}} \\
            }
        }
    }
     \\ \\

    \theutterance \stepcounter{utterance}  
    & & & \multicolumn{2}{p{0.3\linewidth}}{
        \cellcolor[rgb]{0.9,0.9,0.9}{
            \makecell[{{p{\linewidth}}}]{
                \texttt{\tiny{[GM$|$GM]}}
                \texttt{game successful} \\
            }
        }
    }
    & & \\ \\

    \theutterance \stepcounter{utterance}  
    & & & \multicolumn{2}{p{0.3\linewidth}}{
        \cellcolor[rgb]{0.9,0.9,0.9}{
            \makecell[{{p{\linewidth}}}]{
                \texttt{\tiny{[GM$|$GM]}}
                \texttt{end game} \\
            }
        }
    }
    & & \\ \\

\end{supertabular}
}

\end{document}
